% Gantry-Augmentation-Formula-Derivation.tex
% Compile: pdflatex Gantry-Augmentation-Formula-Derivation.tex  (twice for refs)
%
% Sections follow the execution order of build_model() in
% gantry_interconnect_dynamic.py, so the document can be read alongside the code.
%
% Layout: equation(s) -> prose explaining each symbol -> exact code snippet.

\documentclass[11pt,a4paper]{article}

\usepackage{amsmath, amssymb, bm}
\usepackage{xcolor}
\usepackage{listings}
\usepackage{geometry}
\geometry{margin=2.5cm}
\usepackage{parskip}
\usepackage{booktabs}
\usepackage{hyperref}
\hypersetup{colorlinks=true, linkcolor=blue, urlcolor=blue}
\usepackage{caption}
\usepackage{microtype}

% ── Python listing style ──────────────────────────────────────────────────────
\definecolor{codebg}{RGB}{245,245,245}
\definecolor{codeframe}{RGB}{180,180,180}
\definecolor{keywordcolor}{RGB}{0,0,160}
\definecolor{commentcolor}{RGB}{60,110,60}
\definecolor{stringcolor}{RGB}{150,0,0}
\definecolor{numbercolor}{RGB}{100,100,100}

\lstdefinestyle{py}{
  language=Python,
  basicstyle=\ttfamily\footnotesize,
  keywordstyle=\color{keywordcolor}\bfseries,
  commentstyle=\color{commentcolor}\itshape,
  stringstyle=\color{stringcolor},
  showstringspaces=false,
  breaklines=true,
  breakatwhitespace=true,
  frame=single,
  framerule=0.4pt,
  rulecolor=\color{codeframe},
  backgroundcolor=\color{codebg},
  numbers=left,
  numberstyle=\tiny\color{numbercolor},
  numbersep=8pt,
  xleftmargin=14pt,
  captionpos=t,
}

% ── Shorthand ────────────────────────────────────────────────────────────────
\newcommand{\file}[1]{\texttt{\small #1}}
\newcommand{\code}[1]{\texttt{#1}}
\newcommand{\mat}[1]{\bm{#1}}
\newcommand{\norm}[1]{\bar{#1}}

% ── Helper: code block title ─────────────────────────────────────────────────
% Usage: \codetitle{filename}{lines}
\newcommand{\codetitle}[2]{\small\texttt{#1}\quad lines #2}

% =============================================================================
\begin{document}

\begin{center}
  {\LARGE\bfseries Gantry Augmentation: Formula Derivation}\\[6pt]
  {\large Mapping from mathematical description to implementation}\\[4pt]
  {\small Sections follow the execution order of \file{build\_model()} in
  \file{scripts/gantry/gantry\_interconnect\_dynamic.py}}
\end{center}

\tableofcontents
\newpage

% =============================================================================
\section{Normalization}
\label{sec:norm}

Normalization is computed once from training data \emph{before}
\code{build\_model()} is called.
Two distinct normalization objects exist and serve different purposes;
both are used throughout all later sections.

% ─────────────────────────────────────────────────────────────────────────────
\subsection{I/O normalization (deepSI \texttt{norm})}

The encoder and loss function operate on zero-mean, unit-variance signals.
Given the concatenated training inputs $u \in \mathbb{R}^{N \times n_u}$
and outputs $y \in \mathbb{R}^{N \times n_y}$, the statistics are:

\begin{equation}
  \mu_u = \frac{1}{N}\sum_k u_k, \quad
  \sigma_u = \operatorname{std}(u), \quad
  \mu_y = \frac{1}{N}\sum_k y_k, \quad
  \sigma_y = \operatorname{std}(y).
  \label{eq:io_stats}
\end{equation}

The normalized signals used everywhere in training are:

\begin{equation}
  \hat{u}_k = \frac{u_k - \mu_u}{\sigma_u}, \qquad
  \hat{y}_k = \frac{y_k - \mu_y}{\sigma_y}.
  \label{eq:io_norm}
\end{equation}

These statistics are stored on \code{fit\_sys.norm} and applied automatically
by deepSI before any encoder or loss call.

\begin{lstlisting}[style=py,
  title={\codetitle{gantry\_interconnect\_dynamic.py}{166--184}},
  firstnumber=166]
u_all = np.concatenate([t.u for t in train_list])
y_all = np.concatenate([t.y for t in train_list])

fs = 1.0 / train_list[0].dt
P_inv_T = np.linalg.inv(P.numpy().T).astype(DTYPE_NP)  # stage -> logical
x_logical_list = []
for t in train_list:
    pos_logical = (P_inv_T @ t.y.T).T        # (N, 3) stage -> logical
    vel_logical = np.diff(pos_logical, axis=0) * fs  # (N-1, 3)
    vel_logical = np.vstack([vel_logical[:1], vel_logical])  # (N, 3)
    x_logical_list.append(np.hstack([pos_logical, vel_logical]))  # (N, 6)
x_all = np.concatenate(x_logical_list)

x_mean = x_all.mean(axis=0).reshape(NX_PHYS, 1).astype(DTYPE_NP)
std_x  = x_all.std(axis=0).reshape(NX_PHYS, 1).astype(DTYPE_NP) + 1e-8
std_u  = u_all.std(axis=0).reshape(nu, 1).astype(DTYPE_NP) + 1e-8
u_mean = u_all.mean(axis=0).reshape(nu, 1).astype(DTYPE_NP)
ystd   = y_all.std(axis=0).astype(DTYPE_NP) + 1e-8
y0     = y_all.mean(axis=0).astype(DTYPE_NP)
\end{lstlisting}

The training data positions $y \in \mathbb{R}^3$ are measured in
\emph{stage} coordinates (hardware frame).
The FP model operates in \emph{logical} (decoupled) coordinates $q \in \mathbb{R}^3$.
The decoupling matrix $P$ relates them; here it is used as:

\begin{equation}
  q = \underbrace{P^{-\top}}_{\substack{\text{stage} \\ \to\,\text{logical}}}\, y^{\mathrm{stage}},
  \qquad
  \dot{q} \approx \frac{\Delta q}{\Delta t}\quad\text{(finite difference at } f_s\text{)}.
  \label{eq:p_pos}
\end{equation}

This is the only place the P-transform is applied to \emph{outputs}
(positions, used here to build state statistics).
In the physical block (Section~\ref{sec:phys}), $P$ is applied to
\emph{inputs} (forces) in the opposite direction.

Note that $(\mu_x, \sigma_x)$ are derived from
\emph{approximate} logical states reconstructed from measurements via
\eqref{eq:p_pos} and finite-difference velocity.
These are used for state normalization and for the physical matrix
normalization below; they are not directly exposed to deepSI's norm object.

The I/O statistics are then written to \code{fit\_sys.norm} inside
\code{build\_model()}:

\begin{lstlisting}[style=py,
  title={\codetitle{gantry\_interconnect\_dynamic.py}{287--290}},
  firstnumber=287]
fit_sys.norm.u0   = u_mean.flatten()
fit_sys.norm.ustd = std_u.flatten()
fit_sys.norm.y0   = y0
fit_sys.norm.ystd = ystd
\end{lstlisting}

% ─────────────────────────────────────────────────────────────────────────────
\subsection{Physical matrix normalization (\texttt{Cd\_norm})}

The known FP output matrix $C_d \in \mathbb{R}^{n_y \times n_x}$ (from the
MATLAB model) relates physical states to physical outputs:
$y = C_d x + D_d u$.
To use it inside the interconnect, which operates on normalized signals,
$C_d$ is rescaled element-wise so that
$\hat{y} = \bar{C}_d \hat{x} + D_d \hat{u}$
holds in normalized coordinates:

\begin{equation}
  \bar{C}_d[i,j] = C_d[i,j] \cdot \frac{\sigma_{x,j}}{\sigma_{y,i}}.
  \label{eq:cd_norm}
\end{equation}

The derivation: substituting $x = \hat{x} \cdot \sigma_x + \mu_x$ and
$y = \hat{y} \cdot \sigma_y + \mu_y$ into $y = C_d x + D_d u$ and
collecting the normalized terms gives \eqref{eq:cd_norm},
assuming $\mu_y \approx C_d \mu_x + D_d \mu_u$ (satisfied at the training
data operating point).

\begin{lstlisting}[style=py,
  title={\codetitle{gantry\_interconnect\_dynamic.py}{187--188}},
  firstnumber=187]
Cd_norm = Cd.numpy() * std_x.flatten()[None, :] / ystd[:, None]  # (3, 6)
Dd_np   = Dd.numpy()                                               # (3, 3)
\end{lstlisting}

$D_d$ does not require rescaling because it maps normalized input
$\hat{u}$ to normalized output $\hat{y}$ directly: the $\sigma_u / \sigma_y$
factors cancel with the deepSI normalization already applied to $u$ and $y$.

% =============================================================================
\section{Augmented State Partitioning}
\label{sec:states}

The combined model state stacks the physical states from the FP model
and the latent states learned by the ANN:

\begin{equation}
  \hat{x}_k =
  \begin{bmatrix} x_k^{\mathrm{phys}} \\ x_k^{\mathrm{aug}} \end{bmatrix}
  \in \mathbb{R}^{n_{xd}}, \quad
  n_{xd} = n_{x,\mathrm{phys}} + n_{x,\mathrm{ann}}.
  \label{eq:state_partition}
\end{equation}

For the gantry: $n_{x,\mathrm{phys}} = 6$ (positions $q_1, q_2, q_3$ and
velocities $\dot{q}_1, \dot{q}_2, \dot{q}_3$ in logical coordinates) and
$n_{x,\mathrm{ann}} = 2$ (targeting absorber displacement $\delta_a$ and
velocity $\dot{\delta}_a$), giving $n_{xd} = 8$.

Two index arrays select the respective rows from $\hat{x}_k$:

\begin{equation}
  I_{\mathrm{phys}} = \{0,1,2,3,4,5\}, \qquad
  I_{\mathrm{aug}}  = \{6,7\}.
  \label{eq:index_sets}
\end{equation}

\begin{lstlisting}[style=py,
  title={\codetitle{gantry\_interconnect\_dynamic.py}{190, 219--220, 251}},
  firstnumber=190]
PHY_IX = np.arange(NX_PHYS)   # [0,1,2,3,4,5]
\end{lstlisting}

\begin{lstlisting}[style=py,
  title={\codetitle{gantry\_interconnect\_dynamic.py}{219--220, 251}},
  firstnumber=219]
NX_ANN = hp['NX_ANN']
nxd = NX_PHYS + NX_ANN
\end{lstlisting}

\begin{lstlisting}[style=py,
  title={\codetitle{gantry\_interconnect\_dynamic.py}{251}},
  firstnumber=251]
AUG_IX = np.arange(NX_PHYS, nxd)   # augmented state indices [6, 7]
\end{lstlisting}

These index sets are used throughout the interconnect wiring
(Section~\ref{sec:wiring}) as \code{selection\_matrix} and
\code{expansion\_matrix} arguments.

% =============================================================================
\section{Physical State Transition}
\label{sec:phys}

% ─────────────────────────────────────────────────────────────────────────────
\subsection{Continuous-time ODE}

The gantry FP model is an LPV rational system in logical coordinates.
The state is $x = [q^\top,\, \dot{q}^\top]^\top \in \mathbb{R}^6$
with $q = [q_1, q_2, q_3]^\top$ the logical positions.
The scheduling variable is $Y = q_3$ (cross-beam position).

The continuous-time derivative is computed in three steps.

\textbf{Step 1: net logical force.}
Stage-coordinate inputs $u \in \mathbb{R}^3$ are transformed to
logical forces via the decoupling matrix $P$, and the spring-damper
restoring force is subtracted:

\begin{equation}
  f_{\mathrm{net}} = -K q - C \dot{q}
    + \underbrace{P^\top u_{\mathrm{stage}}}_{\substack{\text{stage forces} \\ \to\,\text{logical forces}}},
  \label{eq:fnet}
\end{equation}

where $K \in \mathbb{R}^{3\times 3}$ is the stiffness matrix and
$C \in \mathbb{R}^{3\times 3}$ is the damping matrix (both diagonal for the
gantry in logical coordinates).
Here $P^\top$ transforms stage-coordinate actuator forces into the
decoupled logical force frame, the inverse direction of \eqref{eq:p_pos}.

\textbf{Step 2: LFR rational acceleration.}
The inverse mass matrix is represented as a rational polynomial in $Y$:
$M^{-1}(Y) = N(Y) / d(Y)$, with $N(Y) = N_0 + N_1 Y + N_2 Y^2$
(polynomial matrix, Horner form) and
$d(Y) = m_h(\alpha\gamma - \beta^2 + 2\beta m_h Y + m_h(\alpha - m_h)Y^2)$
(scalar polynomial, determinant of $M(Y)$).
The logical acceleration is:

\begin{equation}
  a = \frac{N(Y)}{d(Y)}\, f_{\mathrm{net}}.
  \label{eq:lfr_accel}
\end{equation}

\textbf{Step 3: LFR signal routing and state derivative.}
The LFR latent signals are formed as $z = [a;\, Ya]$ and
$w = Y z$, capturing the $Y$-dependent coupling.
The full state derivative is then:

\begin{equation}
  \dot{x} = A_x x + B_w w + B_u u_{\mathrm{log}}
           = A_{\mathrm{combined}}\, [x;\, w;\, u_{\mathrm{log}}],
  \label{eq:ode}
\end{equation}

where $A_{\mathrm{combined}}$ fuses $A_x$, $B_w$, $B_u$ into a single
matrix multiplication, avoiding separate matrix-vector products per substep.
All matrices are computed once from physical parameters at construction time
and stored as frozen buffers (not trainable).

\begin{lstlisting}[style=py,
  title={\codetitle{model\_augmentation/fit\_systems/blocks.py}{771--822} -- \texttt{Gantry\_State\_Block.deriv()}},
  firstnumber=771]
def deriv(self, x: Tensor, u: Tensor) -> Tensor:
    # x: (batch, 6, 1) normalised   u: (batch, 3, 1) normalised

    # --- denormalise -------------------------------------------------
    x_phys = x * self.std_x + self.x_mean   # (batch, 6, 1)
    u_phys = u * self.std_u + self.u_mean      # (batch, 3, 1)

    # Work in 2D (batch, n) to match LFR signal-flow convention.
    x2 = x_phys.squeeze(-1)          # (batch, 6)
    u2 = u_phys.squeeze(-1)          # (batch, 3)

    # --- P transform: stage forces -> logical forces -----------------
    u_log = u2 @ self.P_mat.T        # (batch, 3)

    # --- net logical force -------------------------------------------
    # fnet = -K@q - C@qdot + u_logical
    fnet = (-(x2[:, :3] @ self.K_mat.T)
            - (x2[:, 3:] @ self.C_mat.T)
            + u_log)                 # (batch, 3)

    # --- LFR loop solve: a = N(Y)/d(Y) @ fnet -----------------------
    if self.Y_op is not None:
        Y_val = self.Y_op                                  # Python scalar
        a = (self.N_op @ fnet.T).T / self.d_op            # (batch, 3)
    else:
        Y   = x2[:, 2]                                     # (batch,)
        dY  = self.mh * (self.alpha * self.gamma_ - self.beta ** 2
                         + 2 * self.beta * self.mh * Y
                         + self.mh * (self.alpha - self.mh) * Y ** 2)
        Y_r = Y.unsqueeze(0)                               # (1, batch)
        n0f = self.N0 @ fnet.T                             # (3, batch)
        n1f = self.N1 @ fnet.T                             # (3, batch)
        n2f = self.N2 @ fnet.T                             # (3, batch)
        a   = (n0f + Y_r * (n1f + Y_r * n2f)).T / dY[:, None]  # Horner
        Y_val = Y[:, None]                                 # (batch, 1)

    # --- LFR latent signals -----------------------------------------
    z = torch.cat([a, Y_val * a], dim=-1)   # (batch, 6)
    w = Y_val * z                            # (batch, 6)  w = Delta(Y).z

    # --- state update through G (NOT directly from a) ---------------
    combined  = torch.cat([x2, w, u_log], dim=-1)   # (batch, 15)
    xdot_phys = combined @ self.A_combined.T         # (batch, 6)

    # --- renormalise and restore trailing dim -----------------------
    xdot = (xdot_phys / self.std_x_1d).unsqueeze(-1)  # (batch, 6, 1)
    return xdot
\end{lstlisting}

% ─────────────────────────────────────────────────────────────────────────────
\subsection{RK4 discretization with upsampling}

To accurately integrate the continuous ODE at the training sample rate
$T_s$, each sample step is divided into $N_{\mathrm{up}}$ equal substeps
of size $h = T_s / N_{\mathrm{up}}$.
Standard fourth-order Runge-Kutta is applied at each substep:

\begin{align}
  k_1^{(s)} &= h \cdot f\!\left(x^{(s)},\, u\right), \notag \\
  k_2^{(s)} &= h \cdot f\!\left(x^{(s)} + \tfrac{1}{2}k_1^{(s)},\, u\right), \notag \\
  k_3^{(s)} &= h \cdot f\!\left(x^{(s)} + \tfrac{1}{2}k_2^{(s)},\, u\right), \notag \\
  k_4^{(s)} &= h \cdot f\!\left(x^{(s)} + k_3^{(s)},\, u\right), \notag \\
  x^{(s+1)} &= x^{(s)} + \frac{k_1^{(s)} + 2k_2^{(s)} + 2k_3^{(s)} + k_4^{(s)}}{6},
  \label{eq:rk4}
\end{align}

for $s = 0, \ldots, N_{\mathrm{up}}-1$, with $x^{(0)} = x_k$ (normalized state
at sample $k$).
The discrete-time state is $x_{k+1}^{\mathrm{phys}} = x^{(N_{\mathrm{up}})}$.
The input $u$ is held constant across all substeps (zero-order hold).

\begin{lstlisting}[style=py,
  title={\codetitle{model\_augmentation/fit\_systems/blocks.py}{757--769} -- \texttt{Gantry\_State\_Block.nonlinear\_function()}},
  firstnumber=757]
def nonlinear_function(self, z: Tensor):
    assert z.size(1) == self.nx + self.nu
    x = z[:, : self.nx, :]
    u = z[:, self.nx :, :]

    for i in range(self.up_sample):
        k1 = (self.Ts / self.up_sample) * self.deriv(x, u)
        k2 = (self.Ts / self.up_sample) * self.deriv(x + k1 / 2, u)
        k3 = (self.Ts / self.up_sample) * self.deriv(x + k2 / 2, u)
        k4 = (self.Ts / self.up_sample) * self.deriv(x + k3, u)
        x = x + (k1 + 2 * k2 + 2 * k3 + k4) / 6
    return x
\end{lstlisting}

\code{up\_sample} is a hyperparameter (default 2 at $f_s = 4000\,\mathrm{Hz}$),
set in \code{DEFAULT\_HP}. The block receives $z = [\hat{x}^{\mathrm{phys}},\, \hat{u}]$
from the interconnect (see Section~\ref{sec:wiring}).

% =============================================================================
\section{Output Equation}
\label{sec:output}

The model output is the sum of two additive contributions:
one from the frozen physical model and one from the learned augmentation.

% ─────────────────────────────────────────────────────────────────────────────
\subsection{Physical output block (frozen)}

The physical output is a fixed linear map from normalized physical states
and normalized input to normalized output:

\begin{equation}
  \hat{y}_k^{\mathrm{phys}} = \bar{C}_d\, \hat{x}_k^{\mathrm{phys}}
                             + D_d\, \hat{u}_k.
  \label{eq:y_phys}
\end{equation}

$\bar{C}_d$ is the physically derived, normalized C matrix from
\eqref{eq:cd_norm}.
It is stored in a \code{register\_buffer} and is \textbf{not} a trainable
parameter: it is the FP model knowledge locked in place.

\begin{lstlisting}[style=py,
  title={\codetitle{model\_augmentation/fit\_systems/blocks.py}{119--124} -- \texttt{Linear\_Output\_Block.forward()}},
  firstnumber=119]
def forward(self, z: Tensor):
    assert z.size(1) == self.nx + self.nu
    x = z[:, : self.nx, :]
    u = z[:, self.nx :, :]
    w = torch.matmul(self.C, x) + torch.matmul(self.D, u)  # type: ignore
    return w
\end{lstlisting}

Note: \code{self.C} and \code{self.D} are \code{register\_buffer}
(set at construction), not \code{nn.Parameter}.
They move with \code{.to(device)} but receive no gradient.

\begin{lstlisting}[style=py,
  title={\codetitle{gantry\_interconnect\_dynamic.py}{243}},
  firstnumber=243]
out_phys = Linear_Output_Block(C=Cd_norm, D=Dd_np)
\end{lstlisting}

% ─────────────────────────────────────────────────────────────────────────────
\subsection{Augmented output block (learned)}

The augmented output maps the learned latent states to the same output
space, also linearly:

\begin{equation}
  \hat{y}_k^{\mathrm{aug}} = C_{\mathrm{aug}}\, \hat{x}_k^{\mathrm{aug}}
                            + D_{\mathrm{aug}}\, \hat{u}_k,
  \label{eq:y_aug}
\end{equation}

where $C_{\mathrm{aug}} \in \mathbb{R}^{n_y \times n_{x,\mathrm{ann}}}$
and $D_{\mathrm{aug}} \in \mathbb{R}^{n_y \times n_u}$ are
\code{nn.Parameter} and are updated during training.
Here $C_{\mathrm{aug}}$ is stored as \code{nn.Parameter}, so unlike the
physical block it \textbf{does} receive gradients.
The difference in implementation is the only way to see which output
contribution is frozen and which is learned.

\begin{lstlisting}[style=py,
  title={\codetitle{model\_augmentation/fit\_systems/blocks.py}{213--218} -- \texttt{Parameterized\_Linear\_Output\_Block.forward()}},
  firstnumber=213]
def forward(self, z: Tensor):
    assert z.size(1) == self.nx + self.nu
    x = z[:, : self.nx, :]
    u = z[:, self.nx :, :]
    w = torch.matmul(self.C, x) + torch.matmul(self.D, u)
    return w
\end{lstlisting}

The forward pass is identical to \code{Linear\_Output\_Block}, but
\code{self.C} and \code{self.D} are declared as \code{nn.Parameter}
(lines 202--203 in \file{blocks.py}).

\begin{lstlisting}[style=py,
  title={\codetitle{gantry\_interconnect\_dynamic.py}{241--245}},
  firstnumber=241]
C_aug_init = np.zeros((ny, NX_ANN), dtype=DTYPE_NP)
C_aug_init[2, 0] = 1e-2   # HEURISTIC: Y <- delta_a
out_phys = Linear_Output_Block(C=Cd_norm, D=Dd_np)
out_aug  = Parameterized_Linear_Output_Block(
    C=C_aug_init, D=np.zeros((ny, nu), dtype=DTYPE_NP), flag_loss_reg=False)
\end{lstlisting}

$C_{\mathrm{aug}}$ is initialized near zero; only the $Y \leftarrow \delta_a$
entry is set to $10^{-2}$ as a physical prior (the absorber primarily displaces
the Y axis).

% ─────────────────────────────────────────────────────────────────────────────
\subsection{Combined output}

Both output blocks connect \emph{additively} to the single output signal $\hat{y}$:

\begin{equation}
  \hat{y}_k = \hat{y}_k^{\mathrm{phys}} + \hat{y}_k^{\mathrm{aug}}
            = \bar{C}_d\, \hat{x}_k^{\mathrm{phys}}
            + D_d\, \hat{u}_k
            + C_{\mathrm{aug}}\, \hat{x}_k^{\mathrm{aug}}
            + D_{\mathrm{aug}}\, \hat{u}_k.
  \label{eq:y_combined}
\end{equation}

The additive wiring is specified in the interconnect (see
Section~\ref{sec:wiring}, lines 271--275).

% =============================================================================
\section{ANN Augmented State Dynamics}
\label{sec:ann}

The augmented states evolve according to a feedforward ANN
that takes the full combined state and normalized input:

\begin{equation}
  x_{k+1}^{\mathrm{aug}} = \phi_{\mathrm{ann}}\!\left(\hat{x}_k,\, \hat{u}_k\right),
  \quad \phi_{\mathrm{ann}} : \mathbb{R}^{n_{xd}+n_u} \to \mathbb{R}^{n_{x,\mathrm{ann}}}.
  \label{eq:ann_dyn}
\end{equation}

The input dimension is $n_{xd} + n_u = 8 + 3 = 11$.
Taking the \emph{full} state $\hat{x}_k$ (not only $x_k^{\mathrm{aug}}$)
as input allows the ANN to condition the augmented dynamics on the
current physical state, for example coupling the absorber velocity
to the gantry velocity.

The output dimension is $n_{x,\mathrm{ann}} = 2$: only the augmented
state rows are updated by the ANN.
The ANN is initialized with zero output weights so the augmented
dynamics start at zero and grow only as training drives them.

\begin{lstlisting}[style=py,
  title={\codetitle{gantry\_interconnect\_dynamic.py}{251--258}},
  firstnumber=251]
AUG_IX = np.arange(NX_PHYS, nxd)   # augmented state indices [6, 7]
ann_block = Static_ANN_Block(
    nz=nxd + nu, nw=NX_ANN,         # CHANGED: nw=NX_ANN -- ANN outputs only to augmented states
    n_nodes_per_layer=hp['n_nodes_per_layer'],
    n_hidden_layers=hp['n_hidden_layers'],
    net=zero_init_feed_forward_nn,
    activation=_act,
)
\end{lstlisting}

The \code{Static\_ANN\_Block.forward()} evaluates the network on its
flattened input $z \in \mathbb{R}^{n_{xd}+n_u}$ and returns $w \in \mathbb{R}^{n_{x,\mathrm{ann}}}$.
The output is then placed into the augmented state rows of
$\hat{x}_{k+1}$ via an expansion matrix (Section~\ref{sec:wiring}).

% =============================================================================
\section{Interconnect Signal Wiring}
\label{sec:wiring}

The \code{Interconnect} class implements the LFR interconnection matrix
$\mathbf{S}$ from Hoekstra et al.\ (2025, Eq.\ 4).
Each \code{connect\_signals} call specifies one block of $\mathbf{S}$
via a connection method (\texttt{concat} or \texttt{additive}) and
an optional routing matrix.

The five wiring rules that define the full augmented model are:

\begin{lstlisting}[style=py,
  title={\codetitle{gantry\_interconnect\_dynamic.py}{265--275}},
  firstnumber=265]
# CHANGED: route ANN output only to augmented state rows [NX_PHYS:nxd]
ic.connect_block_signals(ann_block, ["x", "u"], [])
ic.connect_signals(ann_block, "xp", "additive", expansion_matrix(AUG_IX, nxd))
ic.connect_signals("x", phy_block, "concat", selection_matrix(PHY_IX, nxd))
ic.connect_block_signals(phy_block, ["u"], [])
ic.connect_signals(phy_block, "xp", "additive", expansion_matrix(PHY_IX, nxd))
# Physical: x_phys + u -> y (additive)
ic.connect_signals("x", out_phys, "concat", selection_matrix(PHY_IX, nxd))
ic.connect_block_signals(out_phys, ["u"], ["y"])
# Augmented: x_aug + u -> y (additive; D_aug init=zeros so u contribution starts at 0)
ic.connect_signals("x", out_aug, "concat", selection_matrix(AUG_IX, nxd))
ic.connect_block_signals(out_aug, ["u"], ["y"])
\end{lstlisting}

Reading each rule against the formulas:

\begin{itemize}
  \item \textbf{Lines 265--266} implement \eqref{eq:ann_dyn}:
    the ANN receives the full state $\hat{x}_k$ and $\hat{u}_k$ as input,
    and its output is added into rows $I_{\mathrm{aug}}$ of $\hat{x}_{k+1}$
    via \code{expansion\_matrix(AUG\_IX, nxd)}.
    This matrix is $E_{\mathrm{aug}} \in \mathbb{R}^{n_{xd} \times n_{x,\mathrm{ann}}}$
    with ones at rows $I_{\mathrm{aug}}$ and zeros elsewhere.

  \item \textbf{Lines 267--269} implement \eqref{eq:ode} in discrete time:
    the physical block receives only rows $I_{\mathrm{phys}}$ of $\hat{x}_k$
    via \code{selection\_matrix(PHY\_IX, nxd)}, plus $\hat{u}_k$,
    and its RK4 output is added into rows $I_{\mathrm{phys}}$ of $\hat{x}_{k+1}$.

  \item \textbf{Lines 270--272} implement the physical output \eqref{eq:y_phys}:
    rows $I_{\mathrm{phys}}$ of $\hat{x}_k$ and $\hat{u}_k$ feed
    \code{out\_phys}, whose result is added to $\hat{y}_k$.

  \item \textbf{Lines 273--275} implement the augmented output \eqref{eq:y_aug}:
    rows $I_{\mathrm{aug}}$ of $\hat{x}_k$ and $\hat{u}_k$ feed
    \code{out\_aug}, whose result is added to $\hat{y}_k$.
\end{itemize}

The \code{Interconnect.init\_forward()} method topologically sorts these
rules into an acyclic computation order before the first forward pass,
ensuring no algebraic loops exist.

% =============================================================================
\section{Encoder Initialization}
\label{sec:encoder}

The encoder maps a window of past inputs and outputs to the current state:

\begin{equation}
  \hat{x}_{k|k} = \psi_\theta\!\left(
    \hat{y}_{k-n_a}^{k},\;
    \hat{u}_{k-n_b}^{k}
  \right),
  \label{eq:encoder}
\end{equation}

where the window length $n_a = n_b = 2 n_{xd} + 1$ follows Jan's standard
(rule: history must span at least twice the state dimension).
With $n_a = n_b$ and \code{na\_right=1}, the window includes the current
step $k$ so the reconstructability map can use $\hat{y}_k$ to estimate $\hat{x}_k$.

The encoder is structured as (Hoekstra 2026, Eq.\ 8):

\begin{equation}
  \hat{x}_{k|k} =
  \underbrace{W_{\psi,y}^b\, \hat{y}_{k-n_a}^{k}
            + W_{\psi,u}^b\, \hat{u}_{k-n_b}^{k}}_{\text{physical states (theory init)}}
  \;\oplus\;
  \underbrace{W_{\psi,y}^a\, \hat{y}_{k-n_a}^{k}
            + W_{\psi,u}^a\, \hat{u}_{k-n_b}^{k}}_{\text{augmented states (random init)}}
  \;+\;
  \underbrace{\tilde{\psi}_{\tilde{\theta}}\!\left(\hat{y}_{k-n_a}^{k},\, \hat{u}_{k-n_b}^{k}\right)}_{\text{shared nonlinear correction}}.
  \label{eq:encoder_full}
\end{equation}

% ─────────────────────────────────────────────────────────────────────────────
\subsection{Physical state rows: reconstructability map}

For a linearized LTI model $(A, B, C, D)$ operating in normalized coordinates,
the $n$-step ahead output can be written as (encoder paper Eq.\ 10):

\begin{equation}
  \begin{bmatrix} y_{k+n} \\ \vdots \\ y_k \end{bmatrix}
  = \mathcal{O}_n\, x_{b,k}
  + \mathcal{T}_n\, \begin{bmatrix} u_{k+n} \\ \vdots \\ u_k \end{bmatrix},
  \label{eq:obs_eqn}
\end{equation}

where $\mathcal{O}_n \in \mathbb{R}^{(n+1)n_y \times n_x}$ is the extended
observability matrix and $\mathcal{T}_n$ is the Toeplitz matrix of Markov
parameters (also written $\Gamma_n$ in the code):

\begin{equation}
  \mathcal{O}_n =
  \begin{bmatrix} CA^n \\ CA^{n-1} \\ \vdots \\ C \end{bmatrix}, \qquad
  \mathcal{T}_n \;\text{(block lower-triangular, entries } CA^{j-i-1}B, D\text{)}.
  \label{eq:On_Tn}
\end{equation}

Inverting \eqref{eq:obs_eqn} for $x_{b,k-n}$ and substituting into
the $n$-step shifted state equation
$x_{b,k} = A^n x_{b,k-n} + r_n u_{k-n}^k$
(encoder paper Eq.\ 14) gives:

\begin{equation}
  x_{b,k} = \underbrace{A^n \mathcal{O}_n^{-1}}_{W_{\psi,y}^b}\, y_{k-n}^{k}
           + \underbrace{\left(-A^n \mathcal{O}_n^{-1} \mathcal{T}_n + r_n\right)}_{W_{\psi,u}^b}\,
             u_{k-n}^k,
  \label{eq:reconstruct}
\end{equation}

where $r_n = [A^{n-1}B, A^{n-2}B, \ldots, B, 0]$ is the input convolution
vector and $\mathcal{O}_n^{-1}$ denotes the Moore-Penrose pseudoinverse.
These are the theory-initialized encoder weights (Hoekstra 2026,
Eqs.\ 16--17).

\begin{lstlisting}[style=py,
  title={\codetitle{model\_augmentation/fit\_systems/pre\_encoder.py}{359--394} -- \texttt{linear\_encoder\_init\_aug.\_\_init\_\_()}},
  firstnumber=359]
# THEORY: Hoekstra 2026 Eq. 10 -- Gamma_n is the Toeplitz matrix T_n
Gamma_n = np.zeros(((n + 1) * ny, (n + 1) * nu))
for i in range(n + 1):
    for j in range(i, n + 1):
        flipped_i = n - i
        flipped_j = n - j
        if i != j:
            Gamma_n[
                flipped_i * ny : (flipped_i + 1) * ny,
                flipped_j * nu : (flipped_j + 1) * nu,
            ] = C @ np.linalg.matrix_power(A, j - i - 1) @ B
        else:
            Gamma_n[
                flipped_i * ny : (flipped_i + 1) * ny,
                flipped_j * nu : (flipped_j + 1) * nu,
            ] = D

# THEORY: Hoekstra 2026 Eq. 10 -- O_n is the observability matrix [C; CA; ...; CA^n]
O_n = np.zeros(((n + 1) * ny, nx))
for i in range(n + 1):
    O_n[i * ny : (i + 1) * ny, :] = C @ np.linalg.matrix_power(A, i)
O_inv = np.linalg.pinv(O_n)

# THEORY: Hoekstra 2026 Eq. 13-14 -- A^n and r_n (input-to-state convolution)
A_n = np.linalg.matrix_power(A, n)
gamma_n = np.zeros((nx, (n + 1) * nu))
for i in range(0, n):
    gamma_n[:, i * nu : (i + 1) * nu] = np.linalg.matrix_power(A, n - i - 1) @ B

# THEORY: Hoekstra 2026 Eq. 17 -- W^b_{psi,y} = A^n O_n^{-1}
self.Wb_psi_y = nn.Parameter(
    torch.tensor(A_n @ O_inv, dtype=torch.float32)
)
# THEORY: Hoekstra 2026 Eq. 16 -- W^b_{psi,u} = -A^n O_n^{-1} T_n + r_n
self.Wb_psi_u = nn.Parameter(
    torch.tensor(-A_n @ O_inv @ Gamma_n + gamma_n, dtype=torch.float32)
)
\end{lstlisting}

% ─────────────────────────────────────────────────────────────────────────────
\subsection{Augmented state rows: random initialization}

The augmented state rows $W_{\psi,y}^a$ and $W_{\psi,u}^a$ have no
physical prior and are initialized randomly (Kaiming uniform):

\begin{equation}
  W_{\psi,y}^a \in \mathbb{R}^{n_{x,\mathrm{ann}} \times (n_a+1)n_y}, \quad
  W_{\psi,u}^a \in \mathbb{R}^{n_{x,\mathrm{ann}} \times (n_b+1)n_u}.
  \label{eq:Wa}
\end{equation}

\begin{lstlisting}[style=py,
  title={\codetitle{model\_augmentation/fit\_systems/pre\_encoder.py}{400--405}},
  firstnumber=400]
wa_y = torch.empty(nx_aug, (n + 1) * ny)
wa_u = torch.empty(nx_aug, (n + 1) * nu)
nn.init.kaiming_uniform_(wa_y)
nn.init.kaiming_uniform_(wa_u)
self.Wa_psi_y = nn.Parameter(wa_y)
self.Wa_psi_u = nn.Parameter(wa_u)
\end{lstlisting}

% ─────────────────────────────────────────────────────────────────────────────
\subsection{Convention fix (D-017)}

The reconstructability map \eqref{eq:reconstruct} was derived for
\emph{pure-scaled} signals $y / \sigma_y$ and $u / \sigma_u$.
The deepSI pipeline feeds \emph{mean-subtracted} signals
$(y - \mu_y)/\sigma_y$ and $(u - \mu_u)/\sigma_u$.
Before applying $W^b$ and $W^a$, the window vectors are shifted back
to pure-scaled convention by adding $\mu_y/\sigma_y$ and $\mu_u/\sigma_u$,
and after applying $W^b$ the physical-state offset $\mu_x/\sigma_x$ is
subtracted to return to the pipeline convention:

\begin{lstlisting}[style=py,
  title={\codetitle{model\_augmentation/fit\_systems/pre\_encoder.py}{444--480} -- \texttt{linear\_encoder\_init\_aug.forward()}},
  firstnumber=444]
def forward(self, uhist, yhist):
    state_has_correct_dimension = len(uhist.size()) > 2
    uhist_mod = uhist.view(uhist.size(0), self.nu * (self.nb + 1), 1)
    yhist_mod = yhist.view(yhist.size(0), self.ny * (self.na + 1), 1)

    # CHANGED: D-055 -- apply convention fix before W^b/W^a (pipeline -> pure-scaled)
    if self.fix_enabled:
        uhist_mod = uhist_mod + self.u_off
        yhist_mod = yhist_mod + self.y_off

    # Physical states from W^b (Eq. 16-17), augmented states from W^a (Eq. 8)
    x_b = self.Wb_psi_u @ uhist_mod + self.Wb_psi_y @ yhist_mod  # (batch, nx, 1)
    x_a = self.Wa_psi_u @ uhist_mod + self.Wa_psi_y @ yhist_mod  # (batch, nx_aug, 1)

    # CHANGED: D-055 -- subtract x_mean/std_x from physical states only
    if self.fix_enabled:
        x_b = x_b - self.x_off

    x = torch.cat([x_b, x_a], dim=1)                             # (batch, nx+nx_aug, 1)
    x = x.view(-1, self.nx + self.nx_aug)

    if self.flag_linear_only:
        return x
    else:
        return x + self.net(
            torch.cat((uhist.view(uhist.size(0), -1), yhist.view(yhist.size(0), -1)), dim=1)
        )
\end{lstlisting}

The buffers \code{u\_off}, \code{y\_off}, \code{x\_off} are computed at
construction time from $\mu_u/\sigma_u$, $\mu_y/\sigma_y$, $\mu_x/\sigma_x$
and tiled to match the full window length.

% ─────────────────────────────────────────────────────────────────────────────
\subsection{Normalized system matrices}

The linearized gantry model is obtained at $T_s$ via
\code{gantry\_linearize\_and\_discretize()}, yielding
$(A_d, B_d, C_d^{dt}, D_d^{dt})$.
These are then normalized so that all encoder weight formulas
\eqref{eq:reconstruct}--\eqref{eq:Wa} operate in the same normalized
space as the interconnect:

\begin{lstlisting}[style=py,
  title={\codetitle{gantry\_interconnect\_dynamic.py}{295--313}},
  firstnumber=295]
Ad, Bd, Cd_dt, Dd_dt = gantry_linearize_and_discretize(dt=TS_NEW)

sys_data_with_x = deepSI.System_data(u=u_all, y=y_all)
sys_data_with_x.x = x_phys_all

Ad_bar, Bd_bar, Cd_bar, Dd_bar = normalize_linear_ss_matrices(
    Ad, Bd, Cd_dt, Dd_dt, sys_data_with_x)
\end{lstlisting}

The normalized matrices $(\bar{A}_d, \bar{B}_d, \bar{C}_d, \bar{D}_d)$
are passed to \code{linear\_encoder\_init\_aug} as the $(A, B, C, D)$
arguments, so $\mathcal{O}_n$, $\mathcal{T}_n$, and $r_n$ are all
computed in normalized state-input-output space.

% =============================================================================
\section{Training Objective}
\label{sec:loss}

Training minimizes the multi-step simulation loss (Hoekstra et al.\ 2025,
Eq.\ 5; Beintema et al.\ 2023).
For a batch of $N$ subsegments of length $T$ starting at randomly chosen
times $k_i$, the truncated loss is:

\begin{equation}
  V_{\mathrm{trunc}}(\theta) =
  \frac{1}{2N(T+1)} \sum_{i=1}^{N} \sum_{\ell=0}^{T-1}
  \left\| \hat{y}_{k_i+\ell|k_i} - y_{k_i+\ell} \right\|_2^2,
  \label{eq:loss}
\end{equation}

where $\hat{y}_{k_i+\ell|k_i}$ is simulated from the encoder-initialized
state $\hat{x}_{k_i|k_i} = \psi_\theta(\hat{y}_{k_i-n_a}^{k_i}, \hat{u}_{k_i-n_b}^{k_i})$
by rolling out the interconnect for $T$ steps:
$(\hat{y}_{k+\ell+1|k}, \hat{x}_{k+\ell+1|k}) = \phi_\theta(\hat{x}_{k+\ell|k}, u_{k+\ell})$.

The parameter vector $\theta = (\theta_{\mathrm{aug}}, \theta_{\mathrm{enc}})$
contains both the augmentation parameters
($C_{\mathrm{aug}}$, $D_{\mathrm{aug}}$, ANN weights) and the encoder
parameters ($W^b$, $W^a$, nonlinear net weights).
All are optimized jointly by Adam.

\begin{lstlisting}[style=py,
  title={\codetitle{model\_augmentation/fit\_systems/interconnect.py}{432--456} -- \texttt{SSE\_Interconnect.loss()}},
  firstnumber=432]
def loss(self, uhist, yhist, ufuture, yfuture, **Loss_kwargs):
    x = self.encoder(uhist, yhist) #initialize Nbatch number of states
    errors = []
    for y, u in zip(torch.transpose(yfuture,0,1), torch.transpose(ufuture,0,1)):
        yhat, x = self.hfn(x, u)
        errors.append(nn.functional.mse_loss(y, yhat))
    loss_MSE = torch.mean(torch.stack(errors))

    has_theta_loss = False
    loss_theta = 0
    for m in self.hfn.connected_blocks:
        if isinstance(m, Parameterized_Linear_State_Block):
            loss_theta = loss_theta + m.param_loss()
            has_theta_loss = True
        elif isinstance(m, Parameterized_Linear_Output_Block):
            loss_theta = loss_theta + m.param_loss()
            has_theta_loss = True
        elif isinstance(m, Parameterized_MSD_State_Block):
            loss_theta = loss_theta + nn.functional.mse_loss(
                m.Lambda * m.params, m.Lambda * m.init_params, reduction="sum")
            has_theta_loss = True
    if has_theta_loss:
        return loss_MSE + loss_theta
    else:
        return loss_MSE
\end{lstlisting}

The inner loop iterates over $T$ future steps (\code{nf} in the code),
computes the MSE at each step, and averages.
For the current setup \code{flag\_loss\_reg=False} on \code{out\_aug},
so \code{param\_loss()} returns zero and the total loss reduces to
$V_{\mathrm{trunc}}$ alone.
The rollout horizon is set by \code{nf = int(0.1\,\mathrm{s} / T_s)}
(five settling time constants of the gantry).

\end{document}
